% !PRINTED 78
\weberpage{78}
%<<CONT
als Funktionen von $u$ betrachtet, verwandte $T$-Funktionen zweiter Ordnung mit den gleichen Periodizitätsfaktoren, und daher sind je drei von ihnen linear abhängig. Es ist also
\[
A\,\vartheta_{00}(u)\,\vartheta_{00}(u+v) + B\,\vartheta_{10}(u)\,\vartheta_{10}(u+v) + C\,\vartheta_{11}(u)\,\vartheta_{11}(u+v) = 0,
\]
woraus für $u=0$, $u=\dfrac{1}{2}$:
\[
\begin{aligned}
A\,\vartheta_{00}\vartheta_{00}(v) + B\,\vartheta_{10}\vartheta_{10}(v) &= 0,\\
A\,\vartheta_{01}\vartheta_{01}(v) + C\,\vartheta_{10}\vartheta_{10}(v) &= 0
\end{aligned}
\]
und daraus:

\weq{11}{\begin{aligned}\vartheta_{10}\vartheta_{10}(v)\,\vartheta_{00}(u)\,\vartheta_{00}(u+v) - \vartheta_{00}\vartheta_{00}(v)\,\vartheta_{10}(u)\,\vartheta_{10}(u+v)\\ {}- \vartheta_{01}\vartheta_{01}(v)\,\vartheta_{11}(u)\,\vartheta_{11}(u+v) = 0.\end{aligned}}

Ebenso sind nun auch, wenn $w$ eine dritte Variable ist, die in der Form
\[
\vartheta_{g_1,g_2}(u+v)\,\vartheta_{g_1,g_2}(u+w), \quad \vartheta_{g_1,g_2}(u)\,\vartheta_{g_1,g_2}(u+v+w)
\]
enthaltenen Produkte, als Funktionen von $u$ betrachtet, verwandte $T$-Funktionen zweiter Ordnung mit den gleichen Periodizitätsfaktoren, so daß zwischen je dreien unter ihnen eine lineare Abhängigkeit besteht. Die Koeffizienten bestimmt man wie oben durch spezielle Werte von $u$. So folgt das Formelsystem:

\weq{12}{\begin{aligned}\vartheta_{00}\vartheta_{00}(v+w)\,\vartheta_{00}(u+v)\,\vartheta_{00}(u+w)\\ {}= \vartheta_{00}(u)\vartheta_{00}(v)\vartheta_{00}(w)\vartheta_{00}(u+v+w)\\ {}- \vartheta_{11}(u)\vartheta_{11}(v)\vartheta_{11}(w)\vartheta_{11}(u+v+w).\end{aligned}}
\weq{13}{\begin{aligned}\vartheta_{01}\vartheta_{01}(v+w)\,\vartheta_{01}(u+v)\,\vartheta_{01}(u+w)\\ {}= \vartheta_{01}(u)\vartheta_{01}(v)\vartheta_{01}(w)\vartheta_{01}(u+v+w)\\ {}+ \vartheta_{11}(u)\vartheta_{11}(v)\vartheta_{11}(w)\vartheta_{11}(u+v+w).\end{aligned}}
\weq{14}{\begin{aligned}\vartheta_{10}\vartheta_{10}(v+w)\,\vartheta_{10}(u+v)\,\vartheta_{10}(u+w)\\ {}= \vartheta_{10}(u)\vartheta_{10}(v)\vartheta_{10}(w)\vartheta_{10}(u+v+w)\\ {}+ \vartheta_{11}(u)\vartheta_{11}(v)\vartheta_{11}(w)\vartheta_{11}(u+v+w).\end{aligned}}

Alle diese Formeln können als spezielle Fälle einer allgemeinen Formel aufgefaßt werden, die \gesperrt{Jacobi} aus den Reihenentwickelungen abgeleitet und zur Begründung der Theorie der elliptischen Funktionen verwandt hat\footnote{Theorie der elliptischen Funktionen, aus den Eigenschaften der Thetareihen abgeleitet. \gesperrt{Jacobis} gesammelte Werke, Bd.\,I, S.\,497.}. Diese Formel läßt sich auch folgendermaßen aus dem Begriffe der Thetafunktionen gewinnen.
